\documentclass[11pt,reqno]{amsart}
\usepackage[T1]{fontenc}
\usepackage{lmodern}
\usepackage{amsmath,amssymb,amsfonts,mathtools}
\usepackage{microtype}
\usepackage{needspace}
\usepackage{booktabs}
\usepackage{hyperref}
\hypersetup{hidelinks,pdfauthor={Y. T. Miah},pdftitle={Images of Multilinear Polynomials with a Standard-Polynomial Unit Specialization},pdfsubject={Multilinear polynomial images of matrix algebras},pdfkeywords={L'vov--Kaplansky conjecture; multilinear polynomials; standard polynomial; matrix image; polynomial identities}}
\urlstyle{same}
\newcommand{\M}{M_n(K)}
\newcommand{\slie}{\mathfrak{sl}_n(K)}
\newcommand{\im}{\operatorname{im}}
\newcommand{\tr}{\operatorname{tr}}
\newcommand{\Id}{I}
\newcommand{\Mat}[1]{M_{#1}(K)}
\newcommand{\dd}{\mathrm d}
\newcommand{\doi}[1]{\href{https://doi.org/#1}{\textup{doi: #1}}}
\theoremstyle{plain}
\newtheorem{theorem}{Theorem}[section]
\newtheorem{lemma}[theorem]{Lemma}
\newtheorem{proposition}[theorem]{Proposition}
\newtheorem{corollary}[theorem]{Corollary}
\theoremstyle{definition}
\newtheorem{definition}[theorem]{Definition}
\newtheorem{remark}[theorem]{Remark}
\numberwithin{equation}{section}
\title[Multilinear images with a standard-polynomial specialization]{Images of Multilinear Polynomials with a Standard-Polynomial Unit Specialization}
\author{Y. T. Miah}
\date{8 November 2026}
\subjclass[2020]{Primary 16R10; Secondary 15A24, 16R30}
\keywords{L'vov--Kaplansky conjecture, multilinear polynomial, standard polynomial, matrix polynomial image, matrix identities, regular differential}
\begin{document}
\begin{abstract}
For $r=3,4,5$, consider an ordinary multilinear polynomial
$f(Y,X_1,\ldots,X_{2r})$ whose specialization at $Y=I$ is a nonzero
scalar multiple of the standard polynomial $s_{2r}$.  Over an
algebraically closed field $K$ of characteristic zero, we prove that
for every $n\ge r+1$ the set of its individual values on $M_n(K)$ is
$\mathfrak{sl}_n(K)$ if its trace vanishes identically, and $M_n(K)$
otherwise.  This includes arbitrary proper perturbations of
$s_7,s_9,s_{11}$.  The proof combines exact standard-polynomial core
constructions, additive size transfer, and a mixed-frame compatibility
lemma with a local lifting criterion of Pan.
\end{abstract}
\maketitle

\section{Introduction and statement}
The L'vov--Kaplansky problem asks whether the set of values of an
ordinary multilinear noncommutative polynomial on a full matrix
algebra is always a vector subspace \cite{KBRY2020}.  We study a
class determined by setting one designated variable equal to the
identity.  Throughout, an image means the set of individual
matrix evaluations, rather than its linear span or Zariski closure.

Let $K$ be an algebraically closed field of characteristic zero.
Write $M_n(K)$ for the algebra of $n\times n$ matrices over $K$,
$\mathfrak{sl}_n(K)$ for its trace-zero subspace, and $\im_n f$
for the image of a polynomial $f$ on $M_n(K)$.  For $d\ge1$ the
unnormalised standard polynomial is
\begin{equation}\label{eq:std}
 s_d(X_1,\ldots,X_d)=\sum_{\sigma\in S_d}\operatorname{sgn}(\sigma)X_{\sigma(1)}\cdots X_{\sigma(d)}.
\end{equation}
A polynomial is \emph{ordinary multilinear} if every monomial
contains each of its named variables exactly once.

The hypothesis includes arbitrary proper perturbations of odd standard
polynomials.  Here a multilinear polynomial is \emph{proper} if
substituting the identity in any one of its variables gives the zero
formal polynomial.  Nonproper perturbations are also permitted.

Our proof uses geometric image methods developed by Tang
\cite[Lemmas~2--3]{Tang2026} and Pan
\cite[\S9.3, Corollary~9.7]{Pan2026}, alongside classical
matrix-identity work \cite{AL1950,Leron1973,Rowen1974,Serre2016}.

\Needspace{17\baselineskip}
\begin{theorem}[Standard-unit restriction in three odd degrees]\label{thm:main}
Let $K$ be algebraically closed of characteristic zero.  Let $r\in\{3,4,5\}$ and $n\ge r+1$, and let
\[
 f(Y,X_1,\ldots,X_{2r})\in K\langle Y,X_1,\ldots,X_{2r}\rangle
\]
be any ordinary multilinear polynomial.  Suppose that, as a formal polynomial,
\begin{equation}\label{eq:unit}
 f(I,X_1,\ldots,X_{2r})=\lambda s_{2r}(X_1,\ldots,X_{2r}),\qquad \lambda\in K^\times.
\end{equation}
Then
\begin{equation}\label{eq:main}
 \im_n f=\begin{cases}
  \mathfrak{sl}_n(K),&\tr f(A)=0\text{ for every }A\in\Mat n^{\,2r+1},\\
  M_n(K),&\tr f(A)\ne0\text{ for some }A\in\Mat n^{\,2r+1}.
 \end{cases}
\end{equation}
These are degrees $7$, $9$ and $11$, with matrix orders $n\ge4$, $n\ge5$ and $n\ge6$, respectively.
\end{theorem}

The standard-polynomial classification \cite[Theorem~1.1]{Miah2026Standard}
gives $\im_n s_{2r}=\mathfrak{sl}_n(K)$ in these ranges, yielding
the trace-zero case by setting $Y=I$.  For a polynomial with nonzero
trace, the main challenge is to find a regular rank-one square-zero
value of $s_{2r}$ and a nonzero cyclic trace gradient at the same
tuple of the remaining inputs.  Our mixed-frame lemma makes these
conditions compatible, including at the sharp degree-$2N$ identity
boundary.  Twelve exact core witnesses and additive size transfer
produce the required regular standard-polynomial values.

\section{Local actual-image lifting}\label{sec:local}


\begin{lemma}[Rank-one nilpotent propagation; cf. Tang]\label{lem:propagate}
Let $n\ge2$ and let $F$ be a polynomial map from an affine space into $V$, where $V$ is either $\mathfrak{sl}_n(K)$ or $M_n(K)$.  Suppose its image is stable under conjugation and nonzero scalar multiplication.  If $F(a)=R$ is nonzero, rank one and square zero, and $\dd F_a$ is surjective onto $V$, then the image of $F$ contains every nonscalar matrix in $V$.  If $V=\mathfrak{sl}_n(K)$ and $0$ also lies in the image, the image is all of $V$.
\end{lemma}
\begin{proof}
Surjectivity of the differential between smooth affine varieties makes $F$ smooth at $a$, and its image contains a Zariski-open neighbourhood $U$ of $R$ (after restricting to a neighbourhood of $a$).  Given a nonscalar $T\in V$, choose a basis with $T_{12}\ne0$.  For
$D(t)=\operatorname{diag}(t^{-1},t,1,\ldots,1)$ the matrix
\[
 t^2D(t)TD(t)^{-1}
\]
extends polynomially to $t=0$ and has value $T_{12}E_{12}$ there.  Conjugation and nonzero scaling identify this value with $R$.  The inverse image of the corresponding open neighbourhood of $R$ on the affine $t$-line contains $0$, hence also some $t\in K^\times$.  Stability under conjugation and scaling then puts $T$ in $\im F$.  In the traceless case the only scalar is $0$.
\end{proof}

\Needspace{9\baselineskip}
\begin{lemma}[Scalar values; cf. Tang]\label{lem:scalar}
If an ordinary multilinear polynomial $f$ has an evaluation of nonzero trace on $M_n(K)$, then $\im_n f$ contains every scalar matrix.
\end{lemma}
\begin{proof}
Choose clock-and-shift matrices $U,V$ with $VU=\zeta UV$, where $\zeta$ is a primitive $n$th root of unity.  The $n^2$ matrices $U^aV^b$ ($0\le a,b<n$) form a basis of $M_n(K)$.  When each input of $f$ is chosen from this basis, all monomials in $f$ yield scalar multiples of the same invertible basis matrix, determined by the sums of the exponents.  Expanding a traceful evaluation multilinearly gives one such basis tuple with nonzero trace.  Its value is a nonzero scalar matrix, because a nontrivial clock-and-shift basis matrix has trace zero.  Scaling one input gives every scalar value, including zero.
\end{proof}

For an ordinary multilinear $f(Y,X_1,\ldots,X_m)$ define its \emph{cyclic $Y$-gradient} $H_f$ by rotating each word until $Y$ is first, then deleting $Y$ and summing its coefficient:
\begin{equation}\label{eq:gradient}
 \tr f(Z,A_1,\ldots,A_m)=\tr\bigl(ZH_f(A_1,\ldots,A_m)\bigr).
\end{equation}
The identity holds in every matrix order, with $H_f$ understood as a formal polynomial.

\begin{proposition}[Singleton unit-deletion lifting; Pan]\label{prop:pan}
Let $f(Y,X)$ be ordinary multilinear and put $g(X)=f(I,X)$.  Suppose $g$ is identically traceless on $M_n(K)$.  If there is one tuple $A$ for which
\begin{equation}\label{eq:pan}
 g(A)\ne0,\quad g(A)^2=0,\quad \operatorname{rank}g(A)=1,\quad
 \im\dd g_A=\mathfrak{sl}_n(K),\quad H_f(A)\ne0,
\end{equation}
then $\im_n f=M_n(K)$.
\end{proposition}
\begin{proof}
This is the singleton-active-slot case of \cite[Corollary~9.7]{Pan2026}.  Indeed, at $(Y,X)=(I,A)$ the $X$-input derivatives span $\mathfrak{sl}_n(K)$ and the $Y$ derivative has nonzero trace by \eqref{eq:gradient}.  Thus the full differential of $f$ is onto $M_n(K)$ at the rank-one square-zero value $g(A)$.  Lemma~\ref{lem:propagate} yields every nonscalar and Lemma~\ref{lem:scalar} yields every scalar.
\end{proof}

\section{Regular standard frames and size transfer}\label{sec:cores}
Fix an even integer $m=2r$, write $k=m-1$ and $N\ge r$.  For $B=(B_1,\ldots,B_k)\in M_N(K)^k$ and $u\in K^N$ define
\[
 S_B=s_k(B_1,\ldots,B_k),\qquad L_B(P)=s_m(B_1,\ldots,B_k,P),
\]
\[
 C_{i,B}(P)=s_k(B_1,\ldots,B_{i-1},P,B_{i+1},\ldots,B_k),
 \quad(1\le i\le k),
\]
and the linear map
\begin{equation}\label{eq:gamma}
\begin{aligned}
 \Gamma_{B,u}:M_N(K)&\longrightarrow M_N(K)\oplus(K^N)^k,\\
 P&\longmapsto\bigl(L_B(P),C_{1,B}(P)u,\ldots,C_{k,B}(P)u\bigr).
\end{aligned}
\end{equation}
\par\Needspace{7\baselineskip}
\begin{samepage}
\noindent Call $N$ \emph{admissible for $s_m$} if some pair $(B,u)$ has
$\det S_B\ne0$ and $\Gamma_{B,u}$ injective.  Such pairs form a
nonempty Zariski-open set whenever one exists.
\end{samepage}

\begin{lemma}[Border-core regularity]\label{lem:core}
If $N$ is admissible for $s_m$, the tuple
\begin{equation}\label{eq:frame}
 A_i=\begin{pmatrix}B_i&0\\0&0\end{pmatrix}\ (1\le i\le k),\qquad
 A_m=\begin{pmatrix}0&u\\0&0\end{pmatrix}\quad \text{in }M_{N+1}(K)
\end{equation}
has $s_m(A)$ nonzero, rank one and square zero.  Moreover, $\dd(s_m)_A$ is onto $\mathfrak{sl}_{N+1}(K)$.
\end{lemma}
\begin{proof}
Only terms with $A_m$ last survive in $s_m(A)$, giving
\begin{equation}\label{eq:val}
 s_m(A)=\begin{pmatrix}0&S_Bu\\0&0\end{pmatrix}.
\end{equation}
The injectivity of $\Gamma_{B,u}$ implies $u\ne0$: if $u=0$, then $I_N$ is in the kernel since $s_m(B,I_N)=0$ for even $m$.  Invertibility of $S_B$ therefore makes \eqref{eq:val} nonzero.

For completeness, consider the annihilator of the full original-input differential under the trace pairing.  Variations of the last input in the upper-right and lower-left rectangles give respectively $v\mapsto S_Bv$ and $c\mapsto -cS_B$.  Since $S_B$ is invertible, every annihilator is block diagonal.  Modulo the scalar identity it can be written $\operatorname{diag}(P,0)$.  The upper-left variation of the last input gives $L_B(Z)$; its trace adjoint is $-L_B$, because the alternating trace in $m+1$ slots is unchanged by cyclic rotation (as $m+1$ is odd) and changes sign when the two tested slots are transposed.  Hence its annihilator equation is $L_B(P)=0$.

A lower-left variation $\left(\begin{smallmatrix}0&0\\c&0\end{smallmatrix}\right)$ of $A_i$ contributes to the upper-left output only when this varied letter immediately follows $A_m$ in a word.  Cyclically rotating the remaining alternating word shows that its trace pairing with $\operatorname{diag}(P,0)$ is $-cC_{i,B}(P)u$.  All remaining original-input variations contribute only to already eliminated border directions or zero.  Thus the complete normalized normal equations are exactly $\Gamma_{B,u}(P)=0$.  Injectivity makes the annihilator modulo scalars zero, as required.  Even standard-polynomial evaluations are traceless by cyclic antisymmetry.
\end{proof}

The additive transfer below uses the standard-polynomial core construction of \cite{Miah2026Standard}.

\begin{lemma}[Fixed-polynomial additive transfer]\label{lem:additive}
For a fixed even standard polynomial $s_m$ and the same free last slot in \eqref{eq:gamma}, admissible core sizes are closed under addition.
\end{lemma}
\begin{proof}
Let $(B,u_B)$ and $(D,u_D)$ be admissible cores of sizes $a,b$.  For a single scalar parameter $t$ set
\[
 C_i(t)=\operatorname{diag}(B_i,tD_i)\ (1\le i\le k),\qquad
 u=\binom{u_B}{u_D}.
\]
The endpoint is $s_k(C(t))=\operatorname{diag}(S_B,t^kS_D)$, invertible for $t\ne0$.  The map $L_{C(t)}:P\mapsto s_m(C_1(t),\ldots,C_k(t),P)$ preserves each of the four matrix-block subspaces of $M_{a+b}(K)$.\par\Needspace{6\baselineskip}
At $t=0$ it acts on the upper mixed rectangle by $W\mapsto S_BW$, and on the lower mixed rectangle by $W\mapsto -WS_B$: all other word positions of the free input carry a vanishing $D$-block factor.  Both mixed-block determinants are therefore nonzero polynomials in $t$.

Choose one $t\in K^\times$ avoiding their roots.  If $\Gamma_{C(t),u}(P)=0$, then $L_{C(t)}(P)=0$ first forces both off-diagonal blocks of $P$ to vanish.  For $P=\operatorname{diag}(P_a,P_b)$, the remaining equations are the two original core systems, with the second block multiplied by nonzero powers $t^k$ or $t^{k-1}$.  Their injectivity gives $P_a=P_b=0$.  The construction preserves the fixed polynomial and its free slot and uses the same parameter $t$ for every block.
\end{proof}

\begin{proposition}[Twelve finite cores and all larger sizes]\label{prop:12}
For $r=3,4,5$, each size $N=r,r+1,\ldots,2r-1$ is admissible for $s_{2r}$.  Consequently every $N\ge r$ is admissible, and every $n\ge r+1$ has a tuple $A\in M_n(K)^{2r}$ satisfying the regularity conditions in Lemma~\ref{lem:core}.
\end{proposition}
\begin{proof}
The finite assertion uses the twelve explicit integer pairs $(B,u)$ in \path{certificates/standard_cores.json}.  The verifier and its modular results are described in Appendix~\ref{app:finite}.  At each of the twelve inputs, the integer endpoint $S_B$ has nonzero determinant modulo $101$ and $103$; $\Gamma_{B,u}$ has rank $N^2$ over each of these finite fields, as computed by a signed-subset recurrence and insertion operator.  In particular some integer maximal minor is nonzero, and both conditions hold over every characteristic-zero field.  Table~\ref{tab:cores} records the exact cases and ranks.  Every integer $N\ge r$ is a sum of integers in $[r,2r-1]$: use the listed cases initially, then subtract $r$ inductively.  Apply Lemma~\ref{lem:additive}, followed by Lemma~\ref{lem:core}.
\end{proof}

\section{Detection on mixed frames}\label{sec:mixed}
A nonzero trace gradient must be detected on the \emph{same} regular family as the standard value.  The following lemma applies even at the first nonidentity order, where degree $2N$ polynomial identities exist.

For the tuple in \eqref{eq:frame} and $T=(t_{ij})\in M_m(K)$ write
\[
 (T\cdot A)_i=\sum_{j=1}^m t_{ij}A_j\qquad(1\le i\le m).
\]
\Needspace{12\baselineskip}
\begin{lemma}[Mixed-standard-frame detection]\label{lem:mixed}
Let $m=2r$, $N\ge r$ and $H\in K\langle X_1,\ldots,X_m\rangle$ be a nonzero ordinary multilinear polynomial.  As the $B_i$, $u$ and $T$ vary, $H(T\cdot A(B,u))$ is not identically zero.  Moreover, it is nonzero for some choice with $T\in\operatorname{GL}_m(K)$; its nonvanishing defines a nonempty open condition on the irreducible parameter space of $B,u,T$.
\end{lemma}
\begin{proof}
Suppose first that identifying any two inputs of $H$ makes the resulting polynomial zero.  Polarizing each equality of two inputs in characteristic zero shows that $H$ changes sign under each transposition.  Hence $H=\alpha s_m$ for some $\alpha\ne0$.  There is a frame of the form \eqref{eq:frame} with $s_m(A)\ne0$: take $B$ to consist, in a suitable order, of $E_{00}$ and the pairs $E_{0j},E_{j0}$ for $1\le j\le r-1$, embedded in the first $N$ coordinates, and take $u=e_0$.  In the alternating evaluation, the loop and $r-1$ two-edge excursions can be ordered in $r!$ ways before the final exit $E_{0,N}$; all have the same sign.  Thus $s_m(A)=\pm r!E_{0,N}\ne0$.

Otherwise identifying some pair of inputs gives a nonzero polynomial $h$ in $m-1$ independent variables, one occurring twice.  It is not a polynomial identity of $M_N(K)$.  If $m<2N$, this follows from the minimal-degree bound for matrix polynomial identities \cite{AL1950}.  At the boundary $m=2N$, polarize the repeated variable in $h$.  The resulting nonzero multilinear polynomial $\widetilde h$ has degree $2N$ and is symmetric in its two polarized positions; nonzeroness follows by specializing those positions back to the same variable, which gives $2h$.  The \emph{entire} degree-$2N$ multilinear identity space is the line $Ks_{2N}$ (the classical sharp fact recorded in \cite[pp.~175--176]{Leron1973}).  If $\widetilde h$ were an identity, it would therefore be alternating in those two positions, contradicting its nonzero symmetry.  Thus $h$ is nonzero at some tuple in $M_N(K)^{m-1}$.

Choose $B_1,\ldots,B_{m-1}$ detecting $h$.  A possibly singular mixing matrix $T_0$ can select those $m-1$ old-block matrices, repeating the chosen identified input in the two specified slots.  Then
\[
 H(T_0\cdot A(B,u))=\operatorname{diag}(h(B),0)\ne0.
\]
Therefore at least one matrix entry of $H(T\cdot A(B,u))$ is a nonzero polynomial in all entries of $B,u,T$.  Its product with $\det T$ remains nonzero in the polynomial ring; an evaluation with both factors nonzero exists because $K$ is infinite.  The same argument proves nonemptiness of the corresponding open set.
\end{proof}

\begin{proposition}[Simultaneous regularity and trace-gradient detection]\label{prop:compat}
If $N$ is admissible for $s_m$ and $H$ is any nonzero ordinary multilinear polynomial of degree $m=2r$, there is a \emph{single} tuple $A'\in M_{N+1}(K)^m$ for which
\begin{equation}\label{eq:triple}
\begin{aligned}
 &s_m(A')\ne0,\qquad s_m(A')^2=0,\qquad
 \operatorname{rank}s_m(A')=1,\\
 &\im\dd(s_m)_{A'}=\mathfrak{sl}_{N+1}(K),\qquad H(A')\ne0.
\end{aligned}
\end{equation}
\end{proposition}
\Needspace{10\baselineskip}
\begin{proof}
Select one nonzero $N^2$-minor of $\Gamma_{B,u}$ and multiply it by $\det S_B$, obtaining a polynomial $\Delta(B,u)$ which is not identically zero.  Lemma~\ref{lem:mixed} gives a nonzero coordinate polynomial $q(B,u,T)$ of $H(T\cdot A(B,u))$.\par\Needspace{8\baselineskip}
The product
\[
 \Delta(B,u)\,(\det T)\,q(B,u,T)
\]
is nonzero in the coordinate ring of an affine space.  Choose $B,u,T$ where it does not vanish, and set $A'=T\cdot A(B,u)$.  Lemma~\ref{lem:core} gives a regular rank-one square-zero value at $A$, and alternation yields
\[
 s_m(T\cdot A)=\det(T)s_m(A),\qquad
 \dd(s_m)_{T\cdot A}\circ T=\det(T)\,\dd(s_m)_A,
\]
where the second equality regards $T$ as the invertible linear change of the full tuple of $X$-inputs.  Hence value and full differential regularity survive the mixing, while $H(A')\ne0$ by construction.
\end{proof}

\section{Proof of the image theorem}\label{sec:mainproof}
\begin{proof}[Proof of Theorem~\ref{thm:main}]
Put $m=2r$.  The standard polynomial $s_m$ has identically zero trace.  If every value of $f$ is traceless, substituting $Y=I$ gives
$\im_n f\supseteq\lambda\im_n s_m=\mathfrak{sl}_n(K)$ by
\cite[Theorem~1.1]{Miah2026Standard}; the reverse inclusion is immediate.

Now suppose that $f$ has a value of nonzero trace.  Its cyclic
$Y$-gradient $H_f$ is a nonzero formal multilinear polynomial.  Since
$m=2r<2n$, the minimal identity-degree bound shows that formal and
evaluated trace vanishing agree in this range.

With $N=n-1\ge r$, Proposition~\ref{prop:12} gives admissibility, and
Proposition~\ref{prop:compat} supplies a tuple $A'$ satisfying
\eqref{eq:triple} for $H_f$.  At $Y=I$, the unit restriction gives
\[
 f(I,A')=\lambda s_m(A'),\qquad
 \dd_Xf\big|_{(I,A')}=\lambda\dd(s_m)_{A'}.
\]
The $Y$-direction has nonzero trace because $H_f(A')\ne0$.
Proposition~\ref{prop:pan}, with $g=\lambda s_m$, yields
$\im_n f=M_n(K)$.
\end{proof}

\begin{corollary}[Proper odd-standard perturbations]\label{cor:proper}
For $r\in\{3,4,5\}$, $n\ge r+1$, $\lambda\ne0$, and any proper ordinary multilinear polynomial $h$ of degree $2r+1$, one has
\[
 \im_n(\lambda s_{2r+1}+h)=M_n(K).
\]
\end{corollary}
\begin{proof}
Put $m=2r$ and let $H$ be the cyclic $Y$-gradient of $\lambda s_{m+1}+h$.  Properness gives $(\lambda s_{m+1}+h)(I,X)=\lambda s_m(X)$ and forces the alternating component of $h$ to vanish: a nonzero alternating component would be a multiple of $s_{m+1}$, whose unit specialization is $s_m\ne0$.  Since $m+1$ is odd, cyclic rotation preserves alternating signs, giving
\begin{equation}\label{eq:properalt}
 \sum_{\pi\in S_m}\operatorname{sgn}(\pi)H(A_{\pi(1)},\ldots,A_{\pi(m)})
 =(m+1)m!\lambda s_m(A).
\end{equation}
At a nonzero standard frame from Proposition~\ref{prop:12}, the right side is nonzero.  Thus at least one permuted tuple detects the gradient, and the evaluated trace is not identically zero.  Theorem~\ref{thm:main} applies.  
\end{proof}

\section*{Use of AI}
OpenAI's GPT-6 Pro contributed substantially to the mathematical investigation, development of proofs, computational verification, literature review, and drafting of this manuscript.  The author is responsible for its mathematical claims and final presentation.

\appendix
\section{Exact finite core checks}\label{app:finite}
The finite part of Proposition~\ref{prop:12} consists of $3+4+5=12$ pairs $(r,N)$, shown in Table~\ref{tab:cores}.  The file \path{certificates/standard_cores.json} lists the twelve integer inputs $B_1,\ldots,B_{2r-1},u$.  The accompanying verifier forms $S_B$ by a signed-subset recurrence and forms the columns of $\Gamma_{B,u}$ by a signed insertion formula, both reduced modulo each of $101$ and $103$.  It tests the nonzero endpoint determinant and full rank $N^2$ and records explicit pivot-row minor residues in \texttt{standard\_core\_checks.json} under \texttt{certificates/}.  As all input entries and the defining polynomials are integers, these nonzero residues imply nonzero integer determinants and a nonzero integral maximal minor.  The argument for arbitrary $N\ge r$ is Lemma~\ref{lem:additive}, not a further finite scan.

\begin{table}[htbp]
\centering
\caption{All certified base sizes; each rank is attained modulo both $101$ and $103$.}\label{tab:cores}
\begin{tabular}{@{}c c c c@{}}\toprule
$r$ & standard polynomial & core sizes $N$ & ranks $\operatorname{rank}\Gamma_{B,u}=N^2$\\\midrule
$3$ & $s_6$ & $3,4,5$ & $9,16,25$\\
$4$ & $s_8$ & $4,5,6,7$ & $16,25,36,49$\\
$5$ & $s_{10}$ & $5,6,7,8,9$ & $25,36,49,64,81$\\\bottomrule
\end{tabular}
\end{table}

To reproduce the finite checks, use Python 3.11 or later and NumPy 1.26 or later (below 3), as specified in \path{requirements.txt}. From the source directory, run
\begin{center}
\texttt{python code/verify\_standard\_cores.py}
\end{center}
The script checks exact coverage of the twelve prescribed pairs, computes the modular minors, raises an error on failure, and regenerates the result JSON.

\Needspace{17\baselineskip}
\bibliographystyle{amsplain}
\bibliography{references}
\end{document}
